\chapter{Conclusions}
\label{ch:conclusions}

\section{What was required and what was delivered}

The project asked for feedback control to hold a photovoltaic array near maximum
power under changing irradiance, with a MATLAB photovoltaic model, an MPPT
model and a control analysis, and a physical model containing the panel and the
DC--DC converter. Table~\ref{tab:delivered} records what was produced against
each of those.

\begin{table}[htbp]
\centering
\caption{Requirements against deliverables.}
\label{tab:delivered}
\begin{tabular}{@{}L{0.26\textwidth}L{0.50\textwidth}L{0.16\textwidth}@{}}
\toprule
\textbf{Required} & \textbf{Delivered} & \textbf{Evidence} \\
\midrule
Photovoltaic model & Five-parameter single-diode model with De Soto translation; parameters identified from five simultaneous datasheet constraints; analytic derivatives; exact maximum power point solver & Ch.~\ref{ch:pvmodel}, Tables~\ref{tab:pvfit},~\ref{tab:pvvalid} \\
MPPT model & Perturb and observe in fixed and variable step form, incremental conductance and extremum seeking, behind a common discrete-time interface & Ch.~\ref{ch:mppt} \\
Control analysis & Averaged converter model, analytic linearisation verified against finite differences, operating envelope sweeps, cascaded design with margins across the envelope and a single-loop baseline for comparison & Ch.~\ref{ch:converter},~\ref{ch:control} \\
Physical model, panel & Custom Simscape component solving the implicit diode equation natively & Ch.~\ref{ch:physical} \\
Physical model, DC--DC & Switched boost converter with real transistor, diode, inductor and capacitors at \SI{20}{\kilo\hertz}, with digital pulse width modulation and computation delay & Ch.~\ref{ch:physical}, Fig.~\ref{fig:slxmodel} \\
Beyond the brief & EN 50530-style dynamic benchmark, 27 automated tests, scripted reproducible pipeline, hardware specification and staged build plan & Ch.~\ref{ch:benchmark},~\ref{ch:software},~\ref{ch:physical} \\
\bottomrule
\end{tabular}
\end{table}

\section{Principal findings}

\subsection{The plant is an operating-point-dependent resonance}

The dynamic resistance of the array appears directly in the damping term of the
converter's characteristic polynomial, equation~\eqref{eq:zeta}. Along the I--V
curve at fixed irradiance the resonant frequency moves by under
\SI{4}{\percent} while the damping ratio changes by a factor of \num{44}, from
\num{0.023} to \num{1.035}. The array is, in effect, the damping element of its
own converter, and it stops damping precisely when it is behaving as a current
source.

\subsection{Regulating the input-side voltage avoids the right half plane zero}

The transfer function from duty ratio to photovoltaic voltage,
equation~\eqref{eq:Gvd}, has a constant numerator. The non-minimum-phase
behaviour that limits boost converter bandwidth appears only in the
duty-to-output path. Choosing the input-side variable as the controlled output
removes a fundamental bandwidth limitation at no cost.

\subsection{A single proportional--integral loop is structurally inadequate}

Against a resonance peak of \SI{43}{\decibel}, a conventionally designed
single-loop controller produces three gain crossovers and its worst-case phase
margin falls to \SI{1.5}{\degree}. The cascade holds above \SI{85}{\degree}
everywhere, because at the inner crossover the plant reduces to $V_{dc}/(Ls)$
independently of operating point. The improvement is structural, not a matter of
retuning.

\subsection{The loop gain on power vanishes at the optimum}

The linearised gain from duty ratio to photovoltaic power equals
$\mathrm{d}P/\mathrm{d}V$, which is zero at the maximum power point and changes
sign across it. No fixed linear controller can regulate to an extremum. This is
why the MPPT layer must estimate a gradient rather than regulate a value, and it
explains the form of every algorithm in Chapter~\ref{ch:mppt} from first
principles.

\subsection{Timescale separation is a derived design rule}

The measured voltage loop settling time of \SI{11.9}{\milli\second} sets the
minimum MPPT sample period. The \SI{20}{\milli\second} used throughout was
derived from that measurement rather than assumed, and the whole three-layer
hierarchy follows from the same reasoning.

\subsection{Extremum seeking is rate-independent, hill climbing is not}

On trapezoidal ramps, extremum seeking held \SI{99.85}{\percent} at every ramp
rate from \SI{25}{} to \SI{200}{\watt\per\square\metre\per\second}, with a worst
instantaneous shortfall of \SI{8}{\watt}. Perturb and observe and incremental
conductance degraded to \SI{99.46}{\percent} with shortfalls above
\SI{116}{\watt}. The advantage is frequency separation between the probe dither
and the disturbance, and it is paid for by a permanent dither loss that makes
extremum seeking the worst performer under constant irradiance.

\subsection{An adaptive rule can introduce a new instability}

The variable-step perturb and observe tracker was the best performer on the step
test and the worst on the ramp test, collapsing to \SI{93.1}{\percent} with a
\SI{705}{\watt} shortfall. Its step size grows with $|\mathrm{d}P/\mathrm{d}V|$,
so once the drift phenomenon points it the wrong way, moving away from the peak
increases the gradient, which increases the step, which moves it further. The
adaptation loop has positive feedback and needs the same stability scrutiny as
the main control loop.

\subsection{The averaged model is trustworthy}

The switched circuit agrees with the averaged model to \SI{0.23}{\percent} in
mean voltage, \SI{0.15}{\percent} in mean current, \SI{1.2}{\percent} in ripple
and \SI{0.005}{\percent} in mean power. Every analytical conclusion in the report
inherits its credibility from this check.

\section{Limitations}

The limitations of the physical model were listed in Chapter~\ref{ch:physical}.
At the level of the study as a whole, three are worth restating. Uniform
illumination is assumed throughout, so partial shading and the multiple local
maxima it creates are outside scope. No sensor noise, quantisation or offset is
modelled, which flatters both hill-climbing algorithms and removes the main
practical distinction between them. And no hardware was built, so every result
is a simulation result, validated against a more detailed simulation rather than
against measurement.

\section{Future work}

\begin{enumerate}[leftmargin=1.6em]
  \item \textbf{A drift-free perturb and observe variant.} Add an irradiance
        change detector that suppresses the direction reversal when a large
        power change is accompanied by a small voltage change. This directly
        addresses the failure of Section~\ref{sec:drift} and is the single most
        valuable extension.
  \item \textbf{Sensor models.} Add quantisation, noise and offset, which should
        reveal the practical difference between perturb and observe and
        incremental conductance that the idealised study could not.
  \item \textbf{Gain scheduling on dynamic resistance.} The outer loop bandwidth
        varies from \SI{15}{} to \SI{223}{\hertz}. Scheduling the outer
        proportional gain on $r_d$, which the tracker already measures from its
        own perturbations, would equalise it.
  \item \textbf{Ripple correlation control.} Use the converter's own
        \SI{20}{\kilo\hertz} switching ripple as the dither, giving extremum
        seeking behaviour with no added excitation and therefore no dither loss.
        This requires the switched model and is a natural use of the physical
        model already built.
  \item \textbf{Partial shading and global search.} Extend the array model to
        per-module irradiance with bypass diodes, and add a global scan layer
        above the local tracker.
  \item \textbf{Hardware realisation.} Execute the staged plan of
        Table~\ref{tab:buildplan}, using the simulated figures as quantitative
        acceptance criteria.
\end{enumerate}

\section{Closing remark}

The value of framing maximum power point tracking as a control problem is that
the framing is predictive rather than decorative. The vanishing loop gain at the
optimum explains why tracking algorithms must estimate gradients. The dynamic
resistance explains why a naive single loop is fragile and where it is fragile.
The measured settling time sets the tracker sample period. The frequency
separation between dither and disturbance explains why extremum seeking is
immune to ramps. In each case a control-theoretic quantity accounts for
behaviour that would otherwise have to be discovered empirically and accepted
without explanation.
